\documentclass[11pt,letterpaper]{article}
\usepackage[margin=.7in,headheight=15pt]{geometry}
\usepackage[T1]{fontenc}
\usepackage{lmodern,amsmath,amssymb,tikz,fancyhdr,parskip,hyperref}
\hypersetup{colorlinks=true,urlcolor=blue}
\pdfinfoomitdate=1\pdftrailerid{}\pdfsuppressptexinfo=15
\pagestyle{fancy}\fancyhf{}
\lhead{\small Week 93 / Squaring a rectangle}\rhead{\small Facilitator guide}
\cfoot{\small Bellingham Math Circle / Week 93 / facilitator / \thepage}
\setlength{\parindent}{0pt}\setlength{\parskip}{6pt}
\begin{document}
\section*{Mathematical overview}
Let a rectangle have \emph{positive} sides $a\le b$ and area $A=ab$. Its equal-area square has a unique positive side $s=\sqrt A$, lying between $a$ and $b$: $a^2\le ab\le b^2$.

\textbf{An equal-area reshaping gives a tighter certificate.} Set
\[m=\frac{a+b}{2},\qquad h=\frac{ab}{m}=\frac{2ab}{a+b}.\]
The new sides are $h$ and $m$. They satisfy
\[hm=ab,\qquad a\le h\le s\le m\le b.\]
For unequal sides, $a<h<s<m<b$. A picture explains the central inequality:
\[m^2-ab=\left(\frac{b-a}{2}\right)^2.\]
The average-side square contains all the rectangle's area plus one small square. Simply replacing \emph{both} sides by their average increases the area unless the old rectangle was already square. Division by $m$ restores the fixed area.

\textbf{The error bracket shrinks, with an exact quantitative rule.} If $d=b-a$, then
\[d_{\rm new}=m-h=\frac{(b-a)^2}{2(a+b)}<\frac d2\quad(a<b).\]
Repeated reshaping preserves $A$ and makes nested brackets for $s$, with widths smaller than $d/2^n$ after $n$ moves. Hence the widths become as small as any prescribed positive tolerance. Near equality, the new gap is proportional to the square of the old gap: nonlinear improvement, not merely a fixed amount subtracted each time.

\textbf{No unequal starting rectangle becomes an exact square in finitely many moves.} The new gap is positive whenever the old gap is positive. A decimal display that appears equal does not prove equality. This holds even when $A$ is a perfect square, such as $36$. Conversely a square remains unchanged.

Children should discover an area comparison and a reliable bracket, then explain why the change helps. Grades 3--5 need rectangular area, halves, division and comparing supplied fractions; the early tile/cutting work can suit some grades 2--3 children. No K--1 packet is forced. Experiments with paper suggest the missing-square identity; exact area checks and the rearrangement establish it. The full general convergence proof is an adult/readiness-dependent continuation.
\newpage
\section*{Preparation, launch and pacing}
Per pair: 24 identical square tiles, squared paper, scissors, a ruler, pencil, whiteboard, and a four-function calculator if desired. Print an extra copy of student page 2 for cutting, retaining an intact reference. Page 1 has a 24-by-20 grid with $7$ mm cells for drawing the tile arrangements; it is not a tray for ordinary larger manipulatives. Build with real tiles on the table, then draw the arrangement on the page. Page 2's cutouts use equal horizontal and vertical scaling, $15$ mm per side unit; rectangle A is $90\times30$ mm and its target square is $60\times60$ mm. Rectangle B is $60\times30$ mm and its target is $45\times45$ mm. Cut diagrams need Actual Size, although the area proof uses labels rather than paper tolerances.

Print an extra copy of student page 4 for Problem 4's movable side cards; retain the original page for recording. The cards are \emph{not} drawings of rectangles at their physical side lengths. The exact fractions are supplied to avoid long division. Keep a calculator optional: reasoning about area and order is the work. Never ask a child to cut a strip $1/2716$ of a unit wide. The page 5 diagrams are conceptual and explicitly not to scale.

\textbf{Launch together.} After movement, let children build several shapes with tiles, including rectangles and a square. Show the concrete action of arranging tiles in straight equal rows and check the area without changing the tile count. Start with Problem 1; let children decide how to organize their collection. Introduce the equal-area-square question after they have shapes to compare.

Before Problem 3, use its complete non-task example with one fixed input, a $1$-by-$3$ rectangle. Its area is $3$. Join its side lengths and halve the combined length to obtain the average $2$. Then find the other new side, $3\div2=3/2$. Show the \emph{whole} output rectangle with both labels $2$ and $3/2$, and check $2\cdot3/2=3$. This full update establishes the purpose: make the same area more square-like and narrow the possible side length of its equal-area square. Averaging alone is not the operation.

Prerequisites: area as rows times columns, equal halves, division, and interpreting fractions as lengths. Adult reading and recording are fine; adults should not choose all arrangements or operate all calculations. One anchored adult can support two pairs. At the three-child upper table rotate builder, checker and recorder. A child who has not yet connected area and multiplication can work on tile rectangles, but is not independently doing the later approximation investigation.

Suggested first visit: five minutes movement, five minutes launch, ten minutes Problem 1, fifteen minutes cuts in Problem 2, five minutes complete update demonstration, and fifteen minutes one or more Problem 3 cases. Problems 4--5 make a second visit. Problems 6--7 deepen the explanation when ready. Finishing every page in an hour is not the goal. Scissor handling, cutout fit, calculator conventions and these timings have not been physically rehearsed. All pages remain unpiloted.
\newpage
\section*{Solutions and optional hints: Problems 1--3}
\textbf{1.} The four whole-tile rectangles, ignoring rotation, are $1\times24$, $2\times12$, $3\times8$ and $4\times6$. The last is closest to square by side difference, side ratio, or the width of its equal-area-square bracket. The equal-area square has side between $4$ and $6$. Completeness: choose the smaller whole side $a$. It cannot exceed $4$, because a $5$-by-$5$ rectangle already needs $25$ tiles; among $a=1,2,3,4$, each divides $24$. The paired side is then forced.

Hints: keep completed rectangles or drawings visible so a rotated copy is recognized; ask what happens after the smaller side reaches $5$. Do not turn the task into a preprinted multiplication table. A child's chosen meaning of ``closest'' can be discussed; for these four pairs the usual meanings agree.

\textbf{2.} A: cut the $6\times2$ rectangle into $4\times2$ and $2\times2$ pieces. Put the first across the bottom of the $4\times4$ square and the second in a top corner. The uncovered square has side $2$ and area $4$.

B: cut the $4\times2$ rectangle into $3\times2$ and $1\times2$ pieces. Put the first across the bottom of the $3\times3$ square, rotate the second to occupy a $2\times1$ part of the top row, leaving a $1\times1$ square. Other dissections are welcome; using all pieces without overlap is essential. Area subtraction checks the uncovered areas $16-12=4$ and $9-8=1$, but the fitting construction shows the required shape, not just its area.

Hint: fit the widest large piece first, then turn the leftover strip. Reserve the general recipe for Problem 6's optional proof discussion.

\textbf{3.} The complete updates and resulting side brackets are:
\[
\begin{array}{c|c|c|c}
\text{old sides}&A&m&h\\\hline
2,6&12&4&3\\
3,6&18&9/2&4\\
2,10&20&6&10/3
\end{array}
\]
The new pairs provide the tightest brackets \emph{among the marked old and new sides}: $[3,4]$, $[4,9/2]$, and $[10/3,6]$. These are not claims that no sharper bounds can ever be found. Direct checks are $3^2<12<4^2$, $4^2<18<(9/2)^2$ and $(10/3)^2<20<6^2$.

Why side lengths enclose the square side: if a square's side were below the shorter rectangle side, its area would be below the rectangle's; above the longer side, its area would be above. The new sides multiply to the same old area and lie between the old sides. That explains the improved certificate. A calculator can supply division, but an unqualified rounded result should not become a claimed exact bound.
\newpage
\section*{Solutions and optional hints: Problems 4--7}
\textbf{4.} The order is B, D, A, C:
\[(2,6)\to(3,4)\to(24/7,7/2)\to(336/97,97/28).\]
Each pair has product $12$. The arithmetic means are respectively $4$, $7/2$ and $97/28$; dividing $12$ by each gives the other side. Card C is the tightest bracket. Its width is
\[\frac{97}{28}-\frac{336}{97}=\frac{9409-9408}{2716}=\frac1{2716}.\]
The earlier widths are $4$, $1$ and $1/14$. The fraction calculation is adult checking material, not a demand to do large handwritten divisions. Squaring shows the square side lies between each positive pair because its product is $12$. Optional hint: compare how close the sides are before verifying the order with the move rule.

\textbf{5.} The decimal bracket is $3.46<s<3.47$. Exact multiplication gives
\[3.46^2=11.9716<12<12.0409=3.47^2.\]
Neither is the exact square side. These are terminating decimal numbers whose squares can be calculated exactly, so the inequalities are certificates; a calculator display of $\sqrt{12}$ alone is not. The printed task also permits finer decimal endpoints one hundredth apart, such as $3.455$ and $3.465$; accept any pair with exact square checks. The displayed pair is unique only if endpoints must be multiples of $0.01$, since squaring positive numbers is strictly increasing.

\textbf{6.} Yes, the average-side square always has at least the rectangle's area, with equality exactly for an original square. Put $m=(a+b)/2$ and $d=(b-a)/2$, so $a=m-d$ and $b=m+d$. Cut the rectangle into an $m$-by-$a$ piece and a $d$-by-$a$ strip. Place the first across the bottom of the $m$-square. Rotate the strip to fill an $a$-wide portion of the $d$-high top row. The remaining corner is a $d$-square. The extra area is $d^2=((b-a)/2)^2$. This is a complete dissection proof for every positive rectangle, including the zero-gap equality case.

\textbf{7.} The bracket can become arbitrarily narrow; an unequal starting rectangle can never become an exact square in a finite number of exact moves. Using Problem 6,
\[m-h=\frac{m^2-A}{m}=\frac{d^2}{m}=\frac{(b-a)^2}{2(a+b)}.\]
Since $b-a<a+b$ for positive sides, the new gap is less than half the old gap. Repeatedly halving a positive fixed width eventually beats any chosen positive tolerance. At every finite step the new gap is still positive if the old one was. The distinction is between approaching the square and reaching it; neither a large number of trials nor a rounded equal-looking display decides this.

Hints: use the little missing square from Problem 6 to describe the area that must be removed from the average-side square; its removal as a strip determines the new gap. The formula is optional for children. A verbal or drawn proof of a uniform shrink factor is substantial mathematical work.
\newpage
\section*{Further investigations, sources and limits}
\textbf{Return visits.} Compare two different equal-area rectangles and ask whether the more square-like one always gives the tighter next bracket. Then begin with area $36$ and sides $4,9$, for which the target side is known to be $6$: the update gives $72/13,13/2$, still unequal. The proof of never reaching a square is therefore not merely an irrational-number issue. Another direction is to set a desired tolerance first and find a starting rectangle that needs several reshaping moves before its bracket is sufficiently narrow. Keep a recoverable pair of side labels instead of a long numerical diary.

The same operation on a long-side estimate $x$ is
\[x_{\rm next}=\frac{x+A/x}{2}.\]
This is Newton's method for $x^2-A=0$. Derivatives are not prerequisites for the rectangle proof. The general Newton algorithm can fail from some starts for other functions; our positive-rectangle proof establishes a global guarantee for this particular problem. Do not transfer the guarantee to arbitrary equations. No historical claim about Babylonian or Greek invention is needed.

\textbf{Connections and novelty.} Week 48 builds inner/outer area certificates and Week 51 propagates guaranteed length ranges. Here a geometric transformation repeatedly improves a side-length certificate while preserving one area. Week 20's harmonic averages and energy questions are different. Atlas GA-28/AP-06 already cover bisection and tangent iteration, including a cubic two-cycle; those activities and examples are excluded. This is a new area-preserving arithmetic/harmonic iteration investigation within an already represented analysis field.

\textbf{Mathematical source.} Tobin A. Driscoll and Richard J. Braun, \emph{Fundamentals of Numerical Computation}, section 4.3, equations (4.3.2), (4.3.7) and Definition 4.3.2:
\url{https://fncbook.com/newton/}.
The source supplies Newton's update, quadratic-convergence context, and its general assumptions. The rectangle construction, missing-square identity, nested bounds, positive-side convergence proof, rational examples and student diagrams are independently derived adaptations. Near the square, the exact gap identity explains quadratic improvement; the finite tests do not establish a general convergence theorem by themselves.

\textbf{Pedagogy consulted.} \emph{Math Circle by the Bay}, Preface printed viii--x, supports deep connected topics, manipulatives, independent work and flexible class-dependent pacing. The grades 3--5 gate and separate later proof visit are our design judgments, not findings from a classroom trial or a source claim that this activity succeeds at those ages.

\textbf{Evidence to collect.} Record whether children preserved area during the full update, understood why two rectangle sides bound a square side, used cards to reduce arithmetic load, and found the cut-and-rearrange proof accessible. Note actual adult interventions and distinguish them from hypotheses about why something was difficult. The student workflow has fresh writer, critic, mathematical review and reviser stages; the guide is a separate untested step. Digital mathematics, byte-identical rebuilding and visual QA do not establish physical fit or classroom readiness. Unpiloted; physical rehearsal untested.
\end{document}
